
En effet, soit @@M0@@ dans @@M1@@. Quitte \`a remplacer @@M2@@ par la cl\^oture alg\'ebrique d'une extension transcendante pure, ce qui est loisible gr\^ace \`a 1.6 (i), on peut supposer qu'il existe @@M3@@ \'el\'ements dans @@M4@@, qui sont alg\'ebriquement ind\'ependants sur tous les @@M5@@. Alors les @@M6@@ sections hyperplanes ``g\'en\'eriques'' successives de @@M7@@, d\'efinies au moyen de ces \'el\'ements, sont toujours des @@M8@@-cycles positifs sp\'eciaux de degr\'e @@M9@@.

Raisonnant par r\'ecurrence sur @@M10@@, il reste donc \`a noter que lorsque @@M11@@, alors @@M12@@, et lorsque @@M13@@, alors @@M14@@. Pour ce dernier point, on peut supposer que @@M15@@ lui-m\^eme provient d'un sch\'ema int\`egre @@M16@@, gr\^ace \`a l'injection @@M17@@ ; mais alors @@M18@@, ce qui est @@M19@@ puisque @@M20@@ est un corps.

\textbf{Corollaire 6.11.} (i) \emph{Soient @@M21@@ deux entiers, @@M22@@ pour @@M23@@ @@M24@@, et @@M25@@ @@M26@@. Soient @@M27@@ un corps alg\'ebriquement clos, et @@M28@@ un @@M29@@-cycle positif sp\'ecial de dimension @@M30@@ et degr\'e @@M31@@, \`a polyn\^ome de Hilbert @@M32@@. Alors pour @@M33@@, on a}
@@M34@@
\emph{De plus, @@M35@@ s'immerge comme sous-sch\'ema de @@M36@@ avec @@M37@@, d\'efini par (au plus) @@M38@@ \'equations de degr\'e @@M39@@.}

(ii) (Chow) \emph{Soit @@M40@@ un sch\'ema noeth\'erien. Lorsque @@M41@@ parcourt les extensions alg\'ebriquement closes des @@M42@@, @@M43@@, les @@M44@@-cycles positifs sp\'eciaux de dimension et degr\'e born\'es forment une ``famille limit\'ee'' (soit avec les @@M45@@ consid\'er\'es comme sous-sch\'emas d'un @@M46@@ fix\'e, soit dans un sens abstrait \'evident).}

